% ECONOMICS_PROBLEM: {"id": "C73-49", "title": "Dynamic games with incomplete information", "jel_code": "C73", "source_id": "OP-0004"}
% !TeX program = xelatex
\documentclass[11pt,a4paper]{article}
\usepackage[margin=23mm]{geometry}
\usepackage{fontspec}
\setmainfont{Latin Modern Roman}
\usepackage{amsmath,amssymb,mathtools}
\usepackage{xcolor,enumitem,booktabs,longtable,array,xurl,hyperref}
\definecolor{muted}{HTML}{53636C}
\hypersetup{colorlinks=true,urlcolor=blue,linkcolor=blue}
\setlength{\parindent}{0pt}
\setlength{\parskip}{5pt}
\newcommand{\R}{\mathbb R}
\newcommand{\E}{\mathbb E}
\newcommand{\N}{\mathbb N}
\newcommand{\ind}{\mathbf 1}
\newcommand{\Var}{\operatorname{Var}}
\newcommand{\Cov}{\operatorname{Cov}}
\newcommand{\supp}{\operatorname{supp}}
\DeclareMathOperator*{\argmin}{arg\,min}
\DeclareMathOperator*{\argmax}{arg\,max}
\newcommand{\entrymeta}[3]{{\sffamily\small\color{muted}\textbf{\texttt{#1}}\quad|\quad #2\par #3\par}}
\newcommand{\Problemref}[1]{\texttt{#1}}
\newcommand{\JELref}[2]{\texttt{#2} (JEL \texttt{#1})}
\begin{document}
\section*{Dynamic games with incomplete information}
\textbf{Problem ID:} \texttt{C73-49}\quad \textbf{JEL:} \texttt{C73}\par
% BEGIN ECONOMICS STATEMENT
\label{problem:OP-0004}
\label{atlas:prob:M4}

\noindent\textbf{Original atlas ID:} M4. \textbf{Formulation:} editorial specification of a research family.

\subsection*{Definitions and assumptions}

Consider a finite-horizon stochastic game at dates $t=0,\ldots,T$, with players $N$, public state $x_t\in X$, private state $s_{it}\in S_i$, and action $a_{it}\in A_i$, all on finite sets. Initial-state and signal distributions, state-transition kernels, observation rules, and utilities are indexed by a structural parameter $\gamma\in\Gamma$. Players discount by a specified $\beta\in(0,1]$. Behavioral strategies depend on each player's information history. Let $\mathcal E(\gamma)$ be the sequential-equilibrium assessments of the resulting finite extensive-form game: strategies are sequentially rational and beliefs are limits of Bayes beliefs induced by completely mixed strategies. Let $P_{\gamma,e}^{\rm obs}$ be the induced observable-data law under assessment $e$. Sampling across independent markets is specified separately.

\subsection*{Formal research question}

For an observed population law $P$, characterize
\[
\Gamma_I(P)=\{\gamma\in\Gamma:\exists e\in\mathcal E(\gamma),\ P_{\gamma,e}^{\rm obs}=P\}.
\]
If a policy $d$ changes specified primitives, define its outcome functional $W(\gamma,e;d)$ and the counterfactual range
\[
\mathcal W_I(d)=\{W(\gamma,e;d):\gamma\in\Gamma_I(P),\ e\in\mathcal E_d(\gamma)\}.
\]
Which restrictions on private-state persistence, signals, transitions, and cross-market selection make these sets informative? Develop algorithms with verified equilibrium residuals and valid approximation bounds for their extrema, together with statistical confidence regions that accommodate sampling uncertainty. A comparative policy effect requires a stated joint selection assumption or an explicitly worst-case difference between counterfactual ranges.

\subsection*{Interpretation and scope}

The finite-horizon formulation ensures that the equilibrium concept and belief consistency are well defined. Infinite-horizon applications require their own equilibrium existence and information-state assumptions; a convenient public Markov state does not automatically summarize histories when private information persists. Hidden-state identification and equilibrium multiplicity can each cause observational equivalence, even when a numerical solver returns one equilibrium reliably. Error in solving an equilibrium must be separated from error in estimating parameters and from uncertainty about selecting equilibria after intervention. A proposed scalable method should therefore report feasibility, incentive residuals, treatment of off-path beliefs, and approximation coverage, rather than claiming that optimization success proves global identification.

\subsection*{Source audit and status}

Ericson--Pakes (1995), Bajari--Benkard--Levin (2007), and Aguirregabiria--Mira (2007) are verified. They concern dynamic industry and Markov-game models under specific state, shock, and equilibrium restrictions. They are important antecedents but do not establish a general solution for arbitrary persistent hidden states and endogenous belief histories. In particular, unobserved transitory payoff shocks are not equivalent to unrestricted evolving private-state information. The atlas title is broader than its citations directly support. The extensive-form model, identified sets, and counterfactual selection treatment above are editorial extensions; the entry remains a methodological agenda, with already solved special models rather than one certified open theorem.

\subsection*{References}

\begin{enumerate}

\item[S1] Richard Ericson and Ariel Pakes (1995). \emph{Markov-Perfect Industry Dynamics: A Framework for Empirical Work}. Review of Economic Studies 62(1), 53--82. \url{https://academic.oup.com/restud/article-abstract/62/1/53/1568000}. \textit{Locator:} abstract; industry equilibrium model. \textit{Establishes:} A dynamic industry framework with entry, exit, and firm-specific uncertainty. \textit{Does not establish:} a solution for arbitrary hidden persistent states and unobserved beliefs.

\item[S2] Patrick Bajari, C. Lanier Benkard, and Jonathan Levin (2007). \emph{Estimating Dynamic Models of Imperfect Competition}. Econometrica 75(5), 1331--1370. \url{https://onlinelibrary.wiley.com/doi/10.1111/j.1468-0262.2007.00796.x}. \textit{Locator:} abstract; two-step estimation procedure. \textit{Establishes:} Estimation under Markov-perfect behavior and specified observable-state restrictions. \textit{Does not establish:} identification of every incomplete-information dynamic game.

\item[S3] Victor Aguirregabiria and Pedro Mira (2007). \emph{Sequential Estimation of Dynamic Discrete Games}. Econometrica 75(1), 1--53. \url{https://onlinelibrary.wiley.com/doi/10.1111/j.1468-0262.2007.00731.x}. \textit{Locator:} abstract; nested pseudo-likelihood method. \textit{Establishes:} A sequential estimator addressing computational and multiplicity issues under stated assumptions. \textit{Does not establish:} generic identification under unrestricted persistent private information.

\end{enumerate}

\subsection*{Common conventions: Source mathematical conventions}

Unless an entry states otherwise, agent and firm sets are finite. State and action spaces are measurable, strategies and outcome maps are measurable, probability laws are specified, and expectations exist for the targets under discussion. Infinite-horizon discount factors lie in $(0,1)$ and discounted objectives are finite. Norms, tolerances and complexity input sizes must be specified for computational comparisons. Constants or primitives that appear locally are inputs, not empirical facts asserted by this catalogue.

\textbf{Identification.} A statistical model $m\in\mathcal M$ implies an observable law $P_m$ and target $\theta(m)$. At law $P$, its identified set is
\[
\Theta(P;\mathcal M)=\{\theta(m):m\in\mathcal M,\ P_m=P\}.
\]
Point identification holds when this set is a singleton, or equivalently when $P_{m_1}=P_{m_2}$ implies $\theta(m_1)=\theta(m_2)$ for all compatible models. Sharp bounds mean bounds attained, or approached when the boundary is not attained, by compatible models. Writing this set is a definition, not a solution: the substantive task is to establish informative restrictions and derive or estimate the set. An unrestricted-confounding model may give only trivial bounds.

\textbf{Inference.} Identification, estimability and valid uncertainty quantification are separate questions. Fix $\alpha\in(0,1)$. A confidence procedure $C_n$ over a declared law class $\mathcal P$ has asymptotic uniform coverage at level $1-\alpha$ when
\[
\liminf_{n\to\infty}\inf_{P\in\mathcal P}\Pr_P\{\theta(P)\in C_n\}\geq1-\alpha.
\]
For set-identified targets the coverage event must be defined separately, for example coverage of the full identified set. If an entry uses identification-indexed models instead of uniquely indexed laws, the same coverage convention is applied over those models. Pointwise limits do not establish uniform validity, and asymptotic validity does not guarantee finite-sample precision. Sample sizes, dependence structures and weak-identification sequences are part of each application.

\textbf{Counterfactuals.} $Y_i(a)$ is a potential outcome under a specified intervention, including its timing, implementation and versions. It is not $Y_i$ conditional on voluntarily choosing $a$. A full intervention vector permits interference; shorthand scalar treatment notation does not silently impose no interference. Assignment, selection, attrition, anticipation and measurement must be specified. A claim about an instrument's local effect is not automatically a population policy effect.

\textbf{Economic equilibrium.} A counterfactual model includes best responses, constraints, feasibility, market clearing, state transitions and expectations consistent with the stated information. When several equilibria exist, either a declared selector is an input or the target is set-valued. Comparisons across policy regimes must use the stated selection convention. Reduced-form relationships are not presumed invariant after an equilibrium-changing intervention.

\textbf{Optimization and welfare.} Optimization is over the stated feasible set. If attainment is not established, $\sup$ or $\inf$ and $\varepsilon$-optimal choices replace an asserted maximizer or minimizer. For example, with value $V=\sup_{a\in\mathcal A}W(a)<\infty$, an $\varepsilon$-optimal set is $\{a\in\mathcal A:W(a)\geq V-\varepsilon\}$ for $\varepsilon>0$. Benefits, costs, risk and health or environmental outcomes can be aggregated only using declared units and conversion weights. Interpersonal utility weights, the treatment of fiscal transfers and foreign profits, discounting, and the behavioral welfare criterion are explicit inputs. A comparative-static derivative additionally requires existence and appropriate differentiability of the selected counterfactual.

\textbf{Sources.} Local labels [S1], [S2], and so forth restart in every question. Locators identify the relevant abstract, model, section or result at the level actually checked; a bibliographic or abstract-level verification is not represented as inspection of an entire proof. Working papers are distinguished from later publications and changing versions. Source summaries are paraphrased. All URL checks and editorial formulations refer to the audit date stated above.
% END ECONOMICS STATEMENT
\end{document}
